\documentclass{article}

\usepackage{iclr2026_conference,times}
\usepackage{amsmath,amssymb}
\usepackage{booktabs}
\usepackage{microtype}
\usepackage{xcolor}
\usepackage{url}
\usepackage{hyperref}

\title{Memory as an Evidence Lifecycle: Context-Isolated Orchestration for Long-Term Agents}

% Anonymous submission. Replace only for the camera-ready version.
\author{Anonymous authors}

\newcommand{\method}{\textsc{EviLife}}
\newcommand{\todo}[1]{\textcolor{red}{[TODO: #1]}}

\begin{document}

\maketitle

\begin{abstract}
Long-term memory for language-model agents is commonly framed as a storage and
retrieval problem: interaction histories are compressed into an external store,
relevant items are retrieved, and the retrieved text is appended to the model
context. We argue that this framing misses a central failure mode. For questions
that require temporal reasoning, state updates, or evidence from multiple
sessions, the relevant records are often retrievable but are not organized into
a representation that a bounded-context agent can reliably read. We introduce
\method, a training-free memory-agent architecture that treats memory use as an
\emph{evidence lifecycle}. Immutable raw memories are exposed through a bounded
navigation catalog; a parent agent decomposes the query into evidence tasks;
read-only subagents inspect disjoint source scopes in isolated context windows;
their full outputs are stored in a result ledger; and an evidence compiler
reconciles provenance, dates, conflicts, exclusions, and coverage into a compact
answer contract. The parent receives references and bounded summaries rather
than raw subagent transcripts. The same agent loop is used for all question
types and receives no gold question-type label, abstention marker, embedding, or
reranker. In a preliminary evaluation on a balanced 180-episode subset of
LongMemEval, our current prototype answers 162 episodes correctly (90.0\%) under
an independent LLM judge. Failure analysis identifies incomplete cross-source
coverage and malformed child reports, rather than lexical retrieval, as the
dominant remaining bottlenecks. These findings motivate a shift from optimizing
memory volume or top-$k$ recall alone toward optimizing evidence organization,
context isolation, and coverage-aware control flow.
\end{abstract}

\section{Introduction}

Language-model agents do not possess memory outside the tokens and state made
available at inference time. As an interaction history grows, an agent must
therefore decide what to preserve, what to retrieve, and how to fit retrieved
information into a finite context window. Long-context models reduce but do not
remove this difficulty: relevant evidence can be diluted by unrelated content,
especially when it occurs far from the query or between distractors
\citep{liu2024lost}. Retrieval-augmented generation reduces the amount of text
presented to the model \citep{lewis2020rag}, but a ranked list of similar chunks
is not yet a coherent account of what happened.

This distinction is consequential in long-term conversational memory. A simple
fact lookup may be answered from one turn. In contrast, a temporal question may
require anchoring a relative expression to the question date, a knowledge-update
question may require selecting the newest valid state, and a multi-session
question may require a complete census across several weakly related sources.
In these cases, increasing $k$ can improve evidence recall while simultaneously
making the reader's task harder. The system has found the evidence but has not
made it readable.

Our experiments expose a second limitation of the standard retriever--reader
pipeline: the agent must recognize \emph{when} a query needs evidence
organization. In an early strict evaluation on all 500 LongMemEval episodes, a
generic agent without benchmark-provided routing priors invoked its evidence
reader on only 6 episodes. It performed well on short, single-source paths but
repeated searches on difficult episodes until exhausting its tool budget. This
diagnostic result suggests that reader quality alone is insufficient; memory
reasoning requires an explicit control loop for task decomposition, source
coverage, and stopping.

We propose \method, an architecture in which the parent agent acts as an
orchestrator rather than a bulk reader. It maintains an explicit task list,
delegates bounded evidence questions to isolated read-only subagents, stores
their full reports outside the parent context, and invokes a separate compiler
to construct a provenance-preserving evidence packet. The architecture is
inspired by tool-using coding agents, where model turns and tool-result turns
alternate in a stateful loop, but it adapts that control flow to the distinct
requirements of persistent conversational memory.

The central hypothesis of this work is:

\begin{quote}
For long-term memory reasoning, the principal systems bottleneck is often not
whether relevant text can be retrieved, but whether distributed evidence can be
organized into a compact, trustworthy, and update-aware state before final
answer generation.
\end{quote}

We make the following contributions:
\begin{itemize}
    \item We formulate long-term memory use as an \emph{evidence lifecycle}
    spanning navigation, scoped reading, result isolation, evidence compilation,
    and final answer projection, rather than as a single top-$k$ retrieval call.
    \item We introduce a context-isolated parent--subagent architecture. Full
    child transcripts remain in sidechains and a result ledger, while the parent
    sees only a result identifier and a bounded summary.
    \item We define provenance- and coverage-aware evidence contracts that
    preserve source references, dates, speakers, state transitions, exclusions,
    conflicts, and unresolved information across compression boundaries.
    \item We evaluate a single no-type-prior agent loop across six LongMemEval
    question types. A preliminary balanced 180-episode experiment reaches
    90.0\% judged accuracy without embeddings or a reranker, and provides an
    auditable taxonomy of remaining evidence-lifecycle failures.
\end{itemize}

The present result is a development-stage finding, not yet a claim of
state-of-the-art performance. The final evaluation will freeze the method and
add full-dataset, cross-model, cross-benchmark, cost, and component-ablation
experiments.

\section{Related Work}

\subsection{Long-context models and retrieval-augmented generation}

Retrieval-augmented generation augments a parametric language model with
non-parametric evidence selected from an external corpus
\citep{lewis2020rag}. Long-context models instead expose more source tokens
directly, but their effective use of information can degrade with position and
interference \citep{liu2024lost}. Both approaches primarily determine
\emph{which text} reaches the model. Our work addresses a subsequent systems
problem: how an agent should divide, validate, reconcile, and compress evidence
when the answer depends on several sources or evolving states.

\subsection{Persistent memory for language-model agents}

Several systems give agents writable memory beyond the active context window.
MemoryBank records and updates user information with an importance- and
time-dependent forgetting mechanism \citep{zhong2023memorybank}. MemGPT frames
context management as virtual memory and lets an agent move information between
memory tiers \citep{packer2023memgpt}. Mem0 extracts and consolidates salient
facts into scalable semantic and graph memories \citep{chhikara2025mem0}, while
A-MEM dynamically links and evolves structured notes
\citep{xu2025amem}. HippoRAG uses a graph and personalized PageRank to support
multi-hop retrieval \citep{gutierrez2024hipporag}.

These methods chiefly innovate in memory representation, curation, or retrieval.
\method{} is complementary: it focuses on the execution path between memory and
the final answer. In particular, it isolates raw reading from parent-agent
planning and makes cross-source evidence compilation a first-class operation.

\subsection{Hierarchical and agent-native memory}

ByteRover proposes an agent-native, file-based Context Tree in which the LLM
curates hierarchical knowledge and progressively escalates retrieval without a
vector database \citep{nguyen2026byterover}. HiGMem similarly uses high-level
event summaries to guide selective reading of underlying turns
\citep{cao2026higmem}. APEX-MEM combines append-only history, temporally grounded
property graphs, and a multi-tool retrieval agent for resolving evolving
information \citep{banerjee2026apex}.

Our navigation catalog is related to hierarchical disclosure, but the proposed
unit of organization differs. Rather than treating a hierarchy or graph as the
final memory interface, we create ephemeral, query-specific evidence tasks and
store each child's complete result outside the parent's context. This design
targets coordination loss, context contamination, and incomplete source
coverage during inference. It can in principle operate over a file hierarchy,
a graph, a lexical index, or another underlying memory store.

\subsection{Long-term conversational memory benchmarks}

LoCoMo evaluates question answering and generation over long multi-session
dialogues \citep{maharana2024locomo}. LongMemEval evaluates information
extraction, temporal reasoning, knowledge updates, multi-session reasoning, and
abstention over 500 timestamped histories \citep{wu2025longmemeval}. These
benchmarks show that long context and conventional retrieval remain below human
performance. We use LongMemEval because its source annotations and heterogeneous
question types allow retrieval coverage to be separated from downstream
evidence-reading failures. Unlike type-aware evaluation pipelines, our runtime
does not receive the benchmark's question type or infer abstention from a
question identifier.

\section{Method}

\subsection{Problem formulation}

Let a conversational memory be a set of immutable source records
$\mathcal{M}=\{m_j\}_{j=1}^{N}$. Each record contains content $x_j$, timestamp
$t_j$, speaker $s_j$, and a stable source reference $p_j$. Given a question
$q$ and question date $t_q$, an agent must produce answer $y$ under an active
context budget $B$:
\begin{equation}
    y \sim f_\theta\!\left(q,t_q,\mathcal{E}_q\right),
    \qquad |\mathcal{E}_q| \leq B,
\end{equation}
where $\mathcal{E}_q$ is not a raw top-$k$ list but a compiled evidence state.
The method must construct $\mathcal{E}_q$ without access to a gold question
type or evidence-session annotation.

We distinguish three notions of completeness. \emph{Retrieval coverage} asks
whether relevant source records appear among candidates. \emph{task-scope
coverage} asks whether a delegated child inspected all sources assigned to it.
\emph{answer coverage} asks whether the combined evidence supports every operand
and boundary in the requested answer. A child can complete its local scope
without establishing global answer coverage.

\subsection{Four-layer memory interface}

\method{} exposes memory through four layers, each retaining references to the
layer below:

\paragraph{Raw memory.}
The raw layer is an append-preserving record of the original interaction. Raw
content is never overwritten by summaries. Corrections create a revision or a
new state event, preserving the evidence needed to audit earlier conclusions.

\paragraph{Navigation catalog.}
A bounded catalog exposes lightweight frontmatter: memory identifiers, source
references, dates, speakers, short descriptions, and structural links. The
catalog is a map of available evidence, not evidence itself. It lets the parent
plan complementary reading tasks without ingesting full records.

\paragraph{Event and state ledger.}
Derived entries organize dated events, entities, relations, and state
transitions while pointing back to raw sources. This layer makes chronology and
supersession explicit but remains revisable when new evidence arrives.

\paragraph{Query-specific evidence results.}
For each question, subagents produce ephemeral evidence results containing
supported candidates, exclusions, missing information, conflicts, source
coverage, and provenance. These results serve the current decision and are not
silently promoted to durable memory.

\subsection{Parent-agent orchestration}

The parent receives $(q,t_q)$ and the bounded catalog. It maintains an explicit
todo list $\mathcal{T}$ and alternates model and tool turns. Its available
operations are deliberately narrow:
\begin{equation}
\textsc{TodoWrite},\quad \textsc{ForkSubagent},\quad
\textsc{CompileEvidence}.
\end{equation}
The parent does not directly read large raw memories in the orchestrator-only
configuration. Instead, it converts unresolved parts of the question into one
or more bounded evidence tasks. Examples include locating all candidate sources
for a predicate, reconstructing a dated update chain, or checking whether two
sources refer to the same event. These are operational tasks inferred from the
question and current evidence state, not benchmark question-type labels.

At turn $k$, the parent state is
\begin{equation}
    z_k = (q,t_q,\mathcal{T}_k,\mathcal{C},\mathcal{R}^{\mathrm{brief}}_k),
\end{equation}
where $\mathcal{C}$ is the navigation catalog and
$\mathcal{R}^{\mathrm{brief}}_k$ contains only child result identifiers and
bounded summaries. The parent chooses to delegate, compile, or answer based on
unresolved todos, source coverage, conflicts, and remaining turn budget.

\subsection{Context-isolated evidence subagents}

A fork request is a tuple
\begin{equation}
    a_i=(d_i, u_i, A_i, P_i, K_i),
\end{equation}
where $d_i$ is a short description, $u_i$ is the evidence objective, $A_i$ is a
tool allowlist, $P_i$ is a set of source or catalog references, and $K_i$ is a
turn budget. Each child receives only the task, required references, and
read-only tools such as \textsc{MemorySearch} and \textsc{MemoryRead}. Recursive
forking and memory mutation are disabled.

The child transcript is written to a separate sidechain. After completion, the
runtime validates the report schema and writes the full result $r_i$ to a
result ledger:
\begin{equation}
    \mathcal{L}[\mathrm{id}_i] \leftarrow
    (u_i,r_i,P_i,\mathrm{coverage}_i,\mathrm{status}_i).
\end{equation}
The parent receives only $(\mathrm{id}_i,\mathrm{summary}_{200}(r_i))$. This
prevents natural-language reports and intermediate tool traces from consuming
the parent's planning context, while preserving full evidence for later
compilation. Tool outputs retain stable source references so compression never
turns a claim into unsupported prose.

\subsection{Evidence compilation}

When sufficient child results are available, the parent calls
\textsc{CompileEvidence} with a set of result identifiers. An isolated compiler
reads the corresponding full records and constructs
\begin{equation}
    \mathcal{E}_q = \Gamma(q,t_q,r_1,\ldots,r_m).
\end{equation}
The compiler cannot call retrieval tools or answer the user directly. It must
produce a bounded evidence contract containing:
\begin{itemize}
    \item answerability and the requested answer operation;
    \item included and excluded claims with source references;
    \item dates, speakers, temporal boundaries, and state transitions;
    \item cross-source coverage, conflicts, unexplored sources, and missing
    information; and
    \item a compact answer projection, such as a fact, list, count, duration,
    latest state, preference, or abstention decision.
\end{itemize}

The runtime validates structural invariants before returning the packet to the
parent. For example, a count must equal the number of included task-level units,
and each included unit must carry provenance. A deterministic fallback preserves
valid child claims when the compiler output is truncated or malformed. The
final paper will separately ablate generic schema validation from any
predicate-specific projection rules present in the development prototype.

\subsection{Progressive disclosure and context compression}

The architecture treats every context boundary as a lossy channel and therefore
defines what must survive it. Raw records are disclosed only after catalog
navigation; child transcripts are replaced by ledger receipts in the parent;
and the compiler emits a final packet bounded by a fixed character or token
budget. Compression preserves user--assistant and tool-call--tool-result pairs,
canonicalizes streamed model output, and keeps reasoning content separate from
answer and tool-input channels. When a packet exceeds its budget, optional
explanations are removed before claims, dates, decisions, or provenance.

This design yields a simple division of responsibility:
\begin{center}
\begin{tabular}{lll}
\toprule
Component & Sees & Responsibility \\
\midrule
Parent agent & question, catalog, result receipts & goals and decisions \\
Evidence child & bounded task and raw sources & scoped discovery and reading \\
Result ledger & complete child reports & lossless intermediate storage \\
Compiler & selected full reports & reconciliation and answer contract \\
\bottomrule
\end{tabular}
\end{center}

\subsection{Runtime contract}

All benchmark question types use the same model loop and tool set. At answer
time, the system receives only the conversation history, source timestamps,
question, and question date. It does not receive \texttt{question\_type}, inspect
question-ID suffixes, use gold evidence annotations, or force a type-specific
reader. The current implementation uses lexical and structural navigation and
does not require embeddings, a vector database, or a learned reranker. This
restriction is a design choice rather than a claim that semantic retrieval is
unnecessary; the orchestration layer is intended to be compatible with stronger
underlying stores.

\section{Experimental Plan}

\paragraph{Current preliminary result.}
On a fixed balanced subset of 180 LongMemEval episodes (30 per runtime-hidden
question type), the current prototype achieves 162/180 judged correct (90.0\%).
Single-session assistant, preference, user, knowledge-update, temporal, and
multi-session accuracy are 100.0\%, 96.7\%, 96.7\%, 90.0\%, 86.7\%, and 70.0\%,
respectively. Only 50.0\% of 332 child reports are structurally valid, and 16 of
18 judged failures contain at least one invalid child report. These diagnostics
support the evidence-lifecycle hypothesis but also show that multi-session
coverage and report reliability remain unresolved.

\paragraph{Required final experiments.}
The submission will report a frozen full-500 LongMemEval run, LoCoMo and at
least one unseen agentic-memory benchmark, multiple backbone models, independent
judges with human agreement checks, token/latency/tool-call cost, and controlled
ablations of the catalog, context isolation, result ledger, compiler, coverage
contract, and deterministic validation. We will compare against full context,
standard lexical and embedding retrieval, a direct parent-only tool agent, and
representative persistent-memory systems. \todo{Run and insert the frozen
publication-grade experiments before making comparative performance claims.}

\section{Limitations}

The present prototype has been iteratively developed on LongMemEval, so the
balanced subset is a regression gate rather than an untouched test set. Its
accuracy is measured by an LLM judge from the same model family as generation,
although calls are independent. The architecture also spends additional model
calls on child reading and compilation. Cross-benchmark generalization,
cross-model robustness, human-judge agreement, and the accuracy--cost frontier
must therefore be established before drawing broad conclusions.

\bibliographystyle{iclr2026_conference}
\bibliography{references}

\end{document}
